% Artigo no formato IEEE (IEEEtran, duas colunas), derivado do ensaio "Quantos bits custa aprender?".
% Compilar: tectonic main.tex   (ou: make artigo, na raiz do repositório)
\documentclass[journal]{IEEEtran}

\usepackage{fontspec}
\setmainfont{texgyretermes}[Extension=.otf,UprightFont=*-regular,BoldFont=*-bold,ItalicFont=*-italic,BoldItalicFont=*-bolditalic]
\usepackage[brazil]{babel}
\usepackage{amsmath,amssymb}
\usepackage{booktabs}
\usepackage{graphicx}
\usepackage{pgfplots}
\pgfplotsset{compat=1.18}
\usepackage{cite}
\usepackage{xcolor}
\usepackage[colorlinks=true,linkcolor=blue!50!black,citecolor=blue!50!black,urlcolor=blue!50!black]{hyperref}

\newcommand{\SNR}{\mathrm{SNR}}
\renewcommand{\IEEEkeywordsname}{Palavras-chave}
% Gerado a partir de experimentos/E12-canal-quantizado-e-ruidoso/resultados/. Provisório até o fim da execução.
\newcommand{\EdozeResumo}{os resultados estão em processamento.}
\newcommand{\EdozeTexto}{Execução em andamento; esta subseção será preenchida com os resultados registrados.}


\begin{document}

\title{Quantos bits custa aprender? Um limite de informação para o treinamento de hardware físico sem backpropagation}

\author{Matheus~Cardoso%
\thanks{M. Cardoso é engenheiro eletricista. Versão 0.2, novembro de 2026. Código, registros e dados em repositório aberto sob licença MIT.}}

\markboth{Pré-publicação, versão 0.2}{Cardoso: Quantos bits custa aprender?}

\maketitle

\begin{abstract}
Hardware analógico, óptico e quântico raramente dispõe de um caminho de volta exato para executar backpropagation e, por isso, estima o gradiente por perturbação através de um canal de medida de banda limitada. Quando cada parâmetro é codificado num tom e separado por detecção síncrona, o limite de Gabor impõe uma janela $T \ge N/B$ para $N$ parâmetros numa banda $B$. Este trabalho propõe que o limite relevante seja dado pela informação exigida pela tarefa, $T \ge I/(B\log_2(1+\SNR))$, em que $I$ é o menor número de bits de realimentação que leva a uma acurácia alvo. Dois experimentos em Fashion-MNIST, com perceptrons de 13.002 a 118.282 parâmetros, testam consequências dessa hipótese em um canal simulado de $r$ projeções aleatórias do gradiente por passo. Com $r$ fixo, a acurácia não diminuiu quando o número de parâmetros cresceu nove vezes (E11; 69,6\% e 74,2\% com $r=8$). Com o canal quantizado e ruidoso (E12), \EdozeResumo{} Os experimentos usam o gradiente exato como fonte das projeções e não contabilizam custo físico de medição; constituem evidência preliminar, e não teste da hipótese em hardware.
\end{abstract}

\begin{IEEEkeywords}
treinamento físico, gradiente por perturbação, capacidade de canal, teoria da informação, computação analógica, computação quântica.
\end{IEEEkeywords}

\section{Introdução}

\IEEEPARstart{O}{} backpropagation obtém o gradiente completo de uma função a um custo que é um múltiplo pequeno do custo de avaliá-la, independentemente do número de parâmetros~\cite{baur1983}. A propriedade depende de armazenar os estados intermediários e percorrer o cálculo em sentido inverso. Em circuitos analógicos, componentes imperfeitos impedem um caminho de volta exato. Em processadores quânticos, a medição destrói o estado e a clonagem é proibida, o que encarece o reaproveitamento de informação entre medidas~\cite{abbas2023}; algumas famílias de circuitos recuperam o escalonamento do backpropagation ao custo de expressividade restrita~\cite{bowles2025}. Em óptica linear e recíproca, por outro lado, o método adjunto fornece o gradiente in situ~\cite{hughes2018,pai2023}, caso discutido na Seção~\ref{sec:fotonica}.

A motivação energética é estimada em mais de dez ordens de grandeza a favor de redes físicas, numa conta que exclui conversão, controle e medição~\cite{tkachenko2025}. Como a medição é a etapa em que o treino acontece, a quantidade de informação que precisa atravessar o canal de medida determina parte relevante desse custo.

As contribuições deste trabalho são:
\begin{enumerate}
\item uma formulação do custo de medida do gradiente como problema de capacidade de canal, a partir da equivalência entre o algoritmo de Jordan e o gradiente multiplexado em frequência (Seção~\ref{sec:formulacao});
\item um experimento (E11) que mede a acurácia em função do número de projeções do gradiente por passo, para redes de tamanhos diferentes;
\item um experimento (E12), com critérios registrados antes da execução, que separa número de medidas, precisão por medida e total de bits;
\item critérios explícitos de refutação e um registro dos resultados negativos que precederam a hipótese.
\end{enumerate}

\section{Trabalhos relacionados}

\paragraph*{Gradiente por perturbação} Métodos de ordem zero estimam o gradiente a partir de avaliações da função. Em otimização convexa, a taxa ótima com duas avaliações por passo cresce com a dimensão~\cite{duchi2015}. O forward gradient usa uma projeção aleatória por passo como estimador sem viés~\cite{baydin2022}, e o MeZO aplica perturbação de ordem zero ao ajuste fino de modelos de linguagem~\cite{malladi2023}. No gradiente multiplexado~\cite{mccaughan2023}, cada parâmetro oscila numa frequência própria e um processo do tipo lock-in separa as contribuições. Oripov et al.~\cite{oripov2025} distinguem duas escalas: o tempo para estimar o gradiente cresce com $N$, mas o tempo até 80\% de acurácia em Fashion-MNIST varia menos de uma ordem de grandeza quando $N$ varia mais de três. Este trabalho procura quantificar em bits o fenômeno observado nessas duas linhas.

\paragraph*{Informação e dimensão do aprendizado} Treinar em subespaços aleatórios define uma dimensão intrínseca da tarefa para uma arquitetura~\cite{li2018}: em redes totalmente conectadas no MNIST, o número de parâmetros variou 24,1 vezes e a dimensão intrínseca, 1,33 vez. A dimensão intrínseca diminui após pré-treino~\cite{aghajanyan2020}. Após um período curto de treino, o gradiente concentra-se num subespaço de dimensão igual ao número de classes, que se mantém ao longo da otimização~\cite{gurari2018}. Medidas em bits do conteúdo aprendido incluem o comprimento de descrição prequencial~\cite{blier2018}, limites PAC-Bayes por compressão~\cite{lotfi2022} e a informação mútua entre pesos e dados~\cite{xu2017}. A quantização do gradiente em poucos bits preserva a convergência em otimização distribuída~\cite{alistarh2017}.

\paragraph*{Treino físico} Além do método adjunto em fotônica~\cite{hughes2018,pai2023}, há treino ciente da física com modelo digital no caminho de volta~\cite{wright2022}, alinhamento por realimentação direta em hardware fotônico~\cite{launay2020,filipovich2022}, propagação de equilíbrio~\cite{scellier2017} e retropropagação por eco hamiltoniano~\cite{lopezpastor2023}. A termodinâmica estocástica relaciona a dissipação mínima à informação adquirida durante o aprendizado~\cite{goldt2017}.

\section{Formulação}
\label{sec:formulacao}

\subsection{Gradiente como coeficiente de Fourier}

No algoritmo de Jordan~\cite{jordan2005}, as perturbações dos parâmetros são preparadas em superposição, a perda é acumulada como fase e a transformada de Fourier quântica lê o gradiente. Para redes clássicas genéricas o algoritmo exige da ordem de $\sqrt{d}$ consultas para $d$ parâmetros~\cite{gilyen2019}, e a carga de dados clássicos em superposição tem custo próprio~\cite{jaques2023}. No gradiente multiplexado, o parâmetro $\theta_i$ recebe a perturbação $\epsilon\cos(2\pi f_i t)$ e a componente da perda na frequência $f_i$ é proporcional a $\partial L/\partial\theta_i$. Nos dois casos, o gradiente é um coeficiente de Fourier da perda perturbada.

\subsection{Limite de leitura e hipótese}

Numa janela de duração $T$, dois tons são resolvidos se estiverem separados por cerca de $1/T$~\cite{gabor1946}. Com $N$ tons numa banda $B$, a leitura completa do gradiente exige
\begin{equation}
T \ge \frac{N}{B}.
\end{equation}
A hipótese deste trabalho substitui $N$ pela informação exigida pela tarefa:
\begin{equation}
\label{eq:hipotese}
T \;\ge\; \frac{I}{B\,\log_2(1+\SNR)},
\end{equation}
em que o denominador é a capacidade do canal gaussiano~\cite{shannon1948} e $I$ é o menor número de bits de realimentação que leva a uma acurácia alvo.

\subsection{Definição operacional de $I$}

Para uma tarefa $\mathcal{D}$, uma arquitetura $A$, uma distribuição de inicialização $p_0$ e um algoritmo $\mathcal{A}$ que recebe realimentação por um canal $C$, define-se
\begin{equation}
\label{eq:I}
I_{\mathcal{A}}(\varepsilon,\delta \mid \mathcal{D},A,p_0) = \min\{\beta : \Pr[\mathrm{acc} \ge \varepsilon] \ge 1-\delta\},
\end{equation}
em que $\beta$ é o total de bits entregues por $C$. Qualquer protocolo que atinge $\varepsilon$ com $\beta$ bits fornece um limite superior para $I_{\mathcal{A}}$. A minimização sobre algoritmos remove a dependência de $\mathcal{A}$; a comparação entre arquiteturas, com $\mathcal{D}$ fixo, estima a parcela atribuível à tarefa.

\subsection{Canal não estacionário e versão quântica}

Durante o treino, os pesos variam e cada tom adquire decaimento. Pela transformada de Laplace, a linha espectral assume forma lorentziana com largura proporcional à taxa de variação, o que prevê menos canais resolvíveis para taxas de aprendizado maiores. A Eq.~\eqref{eq:hipotese} supõe canal estacionário, de modo que a extensão a esse regime é uma questão em aberto. Na versão quântica, o limite de Holevo~\cite{holevo1973} restringe $n$ qubits a no máximo $n$ bits por medida, e o número de execuções do circuito é ao menos $I/n$. A combinação linear de simulações hamiltonianas~\cite{an2023} representa um decaimento como soma de oscilações com peso lorentziano e pode aproximar, em hardware unitário, o decaimento de um gradiente.

\section{Método experimental}

\subsection{Tarefa, modelos e otimização}

Os experimentos usam Fashion-MNIST (60.000 imagens de treino e 10.000 de teste) e perceptrons com duas camadas ocultas de largura $w \in \{16, 32, 64, 128\}$, o que corresponde a $N$ de 13.002 a 118.282 parâmetros. O otimizador é Adam com taxa $10^{-3}$, lotes de 128 e 1.500 passos, com duas sementes.

\subsection{Canal de projeções}

Em cada passo, o gradiente exato $g \in \mathbb{R}^N$ é calculado por backpropagation e projetado em $r$ direções gaussianas de norma unitária, $V \in \mathbb{R}^{N\times r}$. O otimizador recebe $\hat g = V\tilde p$ no lugar de $g$, em que $\tilde p$ é o vetor $p = V^\top g$ após o canal. O estimador não é reescalado por $N/r$. No E11, $\tilde p = p$ em float32. No E12, o canal assume três formas: quantização uniforme estocástica em $2^b$ níveis com escala $s=\max|p|$ transmitida à parte ($rb+32$ bits por passo); ruído gaussiano aditivo com $\SNR = 2^{2b}-1$ ($rb$ bits por passo); e float32 sem degradação, como controle de reprodução do E11.

\subsection{Critérios}

No E11, $r^*(N)$ é o menor $r$ com acurácia a no máximo 2 pontos do backpropagation. No E12, ajusta-se, por $N$ e por canal, a acurácia como função de $\log_2$ do total de bits (critério B1, $R^2 \ge 0{,}90$) e de $\log_2 r$ e $\log_2 b$ separadamente (B2, redução de pelo menos 50\% na soma dos quadrados dos resíduos). B3 exige que, em pelo menos 80\% das configurações, a rede maior não perca mais de 1 ponto. B4 exige que o canal float32 reproduza o E11 com diferença máxima de 0,5 ponto. Os critérios do E12 foram registrados no controle de versão antes da execução.

\section{Resultados}

\subsection{E11: número de projeções}

\begin{figure}[t]
\centering
\begin{tikzpicture}
\begin{axis}[
  width=\columnwidth, height=5.2cm,
  xmode=log, log basis x=2,
  xtick={8,16,32,64,128}, xticklabels={8,16,32,64,128},
  xlabel={$r$ (projeções por passo)}, ylabel={acurácia de teste (\%)},
  ymin=65, ymax=89, grid=major, grid style={gray!20},
  legend style={font=\scriptsize, at={(0.5,-0.32)}, anchor=north, draw=none},
  legend columns=2,
  every axis plot/.append style={thick}]
\addplot[blue, mark=*] coordinates {(8,69.6)(16,74.0)(32,75.8)(64,77.8)(128,78.8)};
\addlegendentry{aleatório, $N$=13.002}
\addplot[blue!50!black, mark=square*] coordinates {(8,74.2)(16,75.9)(32,77.6)(64,79.1)(128,80.7)};
\addlegendentry{aleatório, $N$=118.282}
\addplot[orange, mark=triangle*, dashed] coordinates {(8,74.1)(16,74.0)(32,75.1)(64,76.2)(128,78.0)};
\addlegendentry{adaptativo, $N$=118.282}
\addplot[gray, dotted, domain=8:128, samples=2] {84.5};
\addlegendentry{backprop, $N$=13.002}
\addplot[black, dotted, domain=8:128, samples=2] {87.0};
\addlegendentry{backprop, $N$=118.282}
\end{axis}
\end{tikzpicture}
\caption{E11: acurácia de teste em função do número de projeções por passo (média de duas sementes).}
\label{fig:e11}
\end{figure}

\begin{table}[t]
\centering
\caption{E11: acurácia de teste (\%), modo aleatório}
\label{tab:e11}
\begin{tabular}{rcccc}
\toprule
$N$ & backprop & $r=8$ & $r=32$ & $r=128$ \\
\midrule
13.002 & 84,5 & 69,6 & 75,8 & 78,8 \\
26.506 & 85,6 & 72,6 & 77,3 & 80,1 \\
55.050 & 86,5 & 73,6 & 77,1 & 80,3 \\
118.282 & 87,0 & 74,2 & 77,6 & 80,7 \\
\bottomrule
\end{tabular}
\end{table}

Com $r$ fixo, a acurácia aumentou com $N$ em todas as colunas da Tabela~\ref{tab:e11} (Fig.~\ref{fig:e11}). Com $r=8$, 12.000 projeções levaram a rede de 118.282 parâmetros a 74,2\%, contra 10\% do acaso. Nenhum $r \le 128$ atingiu o critério de 2 pontos, de modo que $r^*$ ficou acima da grade e o expoente de escala não foi definido. Escolher metade das direções a partir das estimativas anteriores reduziu a acurácia (78,0\% contra 80,7\% com $r=128$ na rede maior).

\subsection{E12: medidas, precisão e bits}

\EdozeTexto

\section{Discussão}

\subsection{O que os resultados sustentam}

O E11 é compatível com a hipótese de que o número de canais necessários não cresce com $N$ nessa tarefa, mas não a testa: a acurácia aumentou com $N$, o que também seria esperado se redes maiores fossem apenas mais fáceis de otimizar. O E11 mede número de medidas, não bits; o E12 introduz a precisão por medida como variável. Nenhum dos dois mede a informação mútua efetiva entre as medidas e o gradiente. A concentração do gradiente num subespaço de dimensão próxima ao número de classes~\cite{gurari2018} é compatível com dois resultados do E11: reaproveitar gradientes verdadeiros anteriores aproxima o desempenho do backpropagation com poucas projeções, e reaproveitar estimativas, dominadas por componentes aleatórias, não o faz. Essa leitura sugere que $I$ dependa da estrutura da tarefa, hipótese a testar variando o número de classes.

\subsection{Backpropagation in situ}
\label{sec:fotonica}

O método adjunto em fotônica linear~\cite{hughes2018,pai2023} obtém o gradiente de todas as fases em poucas passagens ópticas, com leitura por intensidade em cada defasador. Nesse esquema o canal de realimentação é largo: a leitura envolve cerca de $N$ detectores e conversores, e o custo desloca-se da duração da janela para o número de detectores e a resolução de conversão. Na formulação da Eq.~\eqref{eq:I}, o backpropagation in situ é um protocolo que fornece um limite superior para $I$ com $\beta \approx N b$ bits por passo. Se $I \ll Nb$, a resolução de leitura pode ser reduzida. O método depende de reciprocidade e de linearidade óptica, e as não linearidades permanecem, em geral, na parte eletrônica.

\section{Limitações e critérios de refutação}

As projeções são calculadas a partir do gradiente exato e não por perturbação física; o estimador não é reescalado; a taxa de aprendizado não foi ajustada por $r$; há uma tarefa e duas sementes; nenhum custo de avaliação de função, janela de integração ou energia foi contabilizado. A hipótese será considerada refutada se $I$ crescer proporcionalmente a $N$ em tarefas realistas, se a Eq.~\eqref{eq:hipotese} errar por ordens de grandeza em bancada, ou se o treino digital permanecer mais barato por joule depois de contabilizadas todas as conversões.

\section{Conclusão}

O custo de treinar hardware sem backpropagation foi formulado como problema de capacidade de canal, com a informação exigida pela tarefa no lugar do número de parâmetros. Em um canal simulado, o número de projeções necessário para uma acurácia dada não aumentou com o tamanho da rede no intervalo testado. Os próximos passos são um desenho fatorial para separar as dependências de $I$ em relação a tarefa, arquitetura, inicialização e algoritmo, a substituição da projeção exata por perturbação física, e a medição em bancada analógica com instrumentação externa.

\appendices
\section{Linhas de pesquisa descartadas}

\begin{table}[h]
\centering
\caption{Hipóteses anteriores e motivo do descarte}
\label{tab:descartadas}
\begin{tabular}{p{0.42\columnwidth}p{0.48\columnwidth}}
\toprule
Hipótese & Resultado \\
\midrule
Base de Fourier fixa nos pesos & sem ganho nas camadas ocultas, como prevê a simetria de permutação \\
Memória em superposição no lugar da atenção & mais energia e menor acerto com chaves semelhantes \\
Custo geométrico da memória & expoente próximo de zero na DRAM até 4~GB \\
Atenção esparsa por conteúdo & $m \propto N^{0,4\text{ a }0,5}$ no GPT-2; mantida como alternativa \\
\bottomrule
\end{tabular}
\end{table}

\section*{Disponibilidade de código e dados}
Scripts, registros completos, sementes e versões das dependências estão no repositório do projeto, com o mapeamento entre cada afirmação e o arquivo de resultado correspondente.

\bibliographystyle{IEEEtran}
\bibliography{referencias}

\end{document}
